\section{A global criterion for three balanced pairs}\label{sec:linked}

We now keep the abelian variety fixed and take relative rank two
with exactly three balanced pairs. After extending the base,
assume it contains a normal closure of \(E\), all the geometric
endomorphisms are defined, and the \(\ell\)-adic algebraic monodromy
group \(G\) is connected. Its rank-two summands under the embeddings
of \(E\) are denoted by \(V_\sigma\).
The regular class function
\[
 B_\sigma(g)=\frac{\operatorname{tr}(g\mid V_\sigma)^2}
                  {\det(g\mid V_\sigma)}
\]
is the character of the adjoint representation plus one. Call
two embeddings linked when these functions agree on \(G\).

\begin{theorem}[Linked projective monodromy]\label{thm:linked}
If two embeddings in different conjugate pairs are linked, then
\(A\) has ordinary reduction at a density-one set of primes
after finite extension.
\end{theorem}

\begin{proof}
At good degree-one primes the relative trace \(\tau\) and determinant
\(\delta\) are algebraic and independent of cohomology, by the
comparison in Section~\ref{sec:coefficients}. Thus equality of two
class functions is equivalent to equality of the corresponding
conjugates of \(\tau^2/\delta\) at all such primes. The reverse
implication follows from Chebotarev: the complement of a class
function's zero set in the Galois image, if nonempty, contains an
open conjugacy-invariant set and is detected in a finite quotient.
The absolute degree-one primes have density one. In particular
linkage is invariant under the rational Galois group on the
embeddings. Polarization gives
\(\bar\tau=p\tau/\delta\) and \(\bar\delta=p^2/\delta\),
so conjugate embeddings are always linked.
On the conjugate pairs, every equivalence class therefore has the
same cardinality, by Galois transitivity.

Take the degree-one primes of the extension in
Theorem~\ref{thm:bounds}, enlarged as above. The trace is nonzero
there. At each split place, the valuations of the projective
trace are \(2v_p(\tau_\sigma)-d_\sigma\); they are odd on
balanced pairs and even on definite pairs. Thus every class is
entirely balanced or entirely definite. A nontrivial linkage
makes the common cardinality greater than one. It divides three
by the number of balanced pairs and hence equals three.
At a balanced component the trace is a unit precisely when the
slopes are \((0,1)\), so equality of the three projective trace
valuations makes all failures simultaneous. By
Corollary~\ref{cor:ranktwo} at most one balanced pair can fail.
Consequently all the selected primes are ordinary.
\end{proof}

Equivalently, the field generated by the good degree-one values
\(\tau^2/\delta\) lies in \(E^+\), and the linkage classes of
conjugate pairs are the fibres of restriction to that field.
Thus a proper projective trace field in \(E^+\) already
implies the conclusion.

\begin{proposition}[Residual geometric endomorphism types]\label{prop:end-frontier}
In the case of three balanced pairs, exclude the potentially CM
varieties and the varieties satisfying Theorem~\ref{thm:linked}.
Up to geometric isogeny the remaining possibilities are:
\begin{enumerate}
\item an absolutely simple variety with full endomorphism algebra
      \(E\) and \(e\ge4\);
\item the square of an absolutely simple threefold whose geometric
      endomorphism field is totally real cubic;
\item an absolutely simple sixfold with a totally indefinite
      quaternion division endomorphism algebra over a totally
      real cubic field.
\end{enumerate}
In these cases the derived monodromy projections on the conjugate
pairs are independent copies of \(\mathrm{SL}_2\), up to central
isogeny.
\end{proposition}

\begin{proof}
Each geometric isotypic component inherits a faithful \(E\)-action.
If there is more than one, the total relative rank two splits as
one plus one, which gives potential CM. Thus write \(A\sim C^n\)
with \(C\) absolutely simple, and let \(F\) and \(t\) be the center
and Schur index of its geometric division endomorphism algebra.

By Faltings's semisimplicity and endomorphism comparison
\cite[Theorems~3--4]{faltings}, over
\(\overline{\Q}_\ell\) the monodromy commutant has matrix blocks of
size \(nt\) indexed by embeddings of \(F\). The corresponding
irreducible monodromy constituents all have the same dimension.
Indeed the \(F\)-component dimension is uniform by the rational
Betti structure. The common constituent dimension is one or two,
as the \(E\)-summands have dimension two. If it is one, the
rational endomorphism algebra admits a commutative semisimple
subalgebra of degree \(2\dim A\) by taking a maximal commutative
\'{e}tale algebra in each matrix factor. This is the potential-CM
case. Such varieties are ordinary at degree-one primes after
adjoining the normal closures of their CM fields of maximal
degree, since each corresponding component is height one.

Otherwise each \(V_\sigma\) is irreducible. The centralizer of
\(E\) in the geometric endomorphism algebra is exactly \(E\),
as its split centralizer has dimension one per \(V_\sigma\).
Therefore \(F\subset E\) and \([E:F]=nt\).
All embeddings above an embedding of \(F\) give isomorphic
representations, by the above commutant decomposition. In the
absence of a cross-pair link, a fibre has size at most two,
and a fibre of size two consists of conjugates.
If the size is two, \(F=E^+\). Equal determinant valuations
at the two embeddings above each place then force every
conjugate pair to be balanced. Hence \(e=3\).
Here \(nt=2\), giving either the real-cubic square or a quaternion
division endomorphism algebra. If the latter is ramified at a
real place, positivity forces the Rosati involution there to be
canonical quaternion conjugation, of symplectic type. An
orthogonal involution would fix a nonzero pure quaternion and
therefore violate positivity. After scalar extension the relevant
polarization on the two-dimensional matrix module and its
two-dimensional multiplicity space is the tensor of an alternating
and a symmetric form. The derived monodromy would then act through
\(\mathrm{SO}_2\) on the multiplicity space, contradicting its
irreducibility since the connected center acts by scalars.
Thus the quaternion algebra is totally indefinite.

A fibre of size one gives \(n=t=1\) and full endomorphism
algebra \(E\). If \(e=3\), all pairs would be balanced and the
relative determinant would have type \((1,1)\) at every embedding.
It is therefore trivial on the Hodge circle, and its kernel is
defined over \(\Q\), so it is trivial on the Hodge group.
Polarization and the self-duality of a two-dimensional representation
of determinant one make the conjugate summands isomorphic as Hodge
group modules. The complex Hodge commutant then contains
off-diagonal maps, whereas the commutant of the rational Hodge
structure equals the algebraic endomorphism algebra of the
abelian variety after tensoring with \(\Q\). This excludes \(e=3\)
in the full-field case.

For any of the remaining types, each derived projection is
\(\mathrm{SL}_2\) since the monodromy is connected reductive and
the two-dimensional summand is irreducible. A nontrivial
projection of the semisimple derived Lie algebra onto
\(\mathfrak{sl}_2\) uses one simple ideal. Two projections
using the same ideal give the same projective trace function,
as they are both the standard two-dimensional action, up to
conjugacy, and scalars cancel. The absence of a cross-pair link
therefore gives distinct factors for distinct pairs.
\end{proof}

In particular, a larger endomorphism algebra with CM center
forces ordinarity in this setting. The remaining real-cubic
square case already contains the ordinary-reduction question for
real-multiplication threefolds over arbitrary number fields.
For such a polarized threefold \(C\) with real field \(F\) fixed
by Rosati, the product \(C^2\) carries the diagonal \(F\)-action
and the standard rotation giving \(F(i)\). The three real
embeddings of \(F\) have equal complementary Hodge dimensions,
so this action has exactly three balanced pairs. Ordinarity of
\(C^2\) and of \(C\) are equivalent. Fité's cubic result
\cite[Theorem~3]{fite} assumes the real field is not contained
in a field of definition of its action, and the general theorem of
Wang and Zhang \cite[Theorem~1.3]{wangzhang} assumes an
auxiliary splitting pattern. These hypotheses need not hold
in the split frontier above.
